\section{準備}
本節では，本研究で扱う言語の構文および宣言的な型付け規則を定義する．型付け規則の表記は\cite{DBLP:conf/popl/DolanM17}に倣う．
\subsection{構文}
型$\tau$はその\textbf{極性}により，正負が割り当てられる．正の型および負の型の構文を次のように帰納的に定める．
\begin{definition} [型の構文]
正の型$P$， 負の型$N$の構文を，以下のように帰納的に定義する．
\nonterms{P}
\nonterms{N}
\end{definition}
項および型の構文は次のように定義する．項構文は PCF をレコード及びヴァリアントにより拡張した体系である PCF+\cite{gunter92}にパターンマッチやADTを加えたものである．
\label{subsec:syntax}
\begin{definition}[項の構文]
項$t$およびパターン$p$を，以下のように帰納的に定義する．
\nonterms{t}
\nonterms{p}
\end{definition}

次に，宣言的な型付け規則を定める．まず，単相的文脈$\Gamma$およびtypingを次のように定めよう．
\begin{definition}[単相的文脈]
単相的文脈$\Gamma$とは，変数から型への有限写像である:
\nonterms{D}
\end{definition}
\begin{definition}[Typing]
Typing $\mathcal{T}$とは，単相的文脈と型の組である:
\nonterms{Tau}
\end{definition}
\begin{definition}[多相的文脈]
多相的文脈$\Pi$はlet-boundな変数からtypingへの有限写像である:
\nonterms{ctx}
\end{definition}
このもとで，項の型付けを以下のように定める．
\begin{figure}[H]
\drules[T] {$\Pi \Vdash t : \tau$} {} {Var-Mono, Var-Poly, Abs, App, Let, Cons, Proj, Variant, Sub, Rec-Mono}
\caption{項の型付け規則(一部抜粋)}
\end{figure}